\chapter{Criterio de salida pública de claims}

\section{Salida pública}
\begin{definition}
Se llama salida pública de claims al paso curatorial y metodológico mediante el cual un subconjunto de claims ya suficientemente estabilizados dentro de Tomo III puede proyectarse hacia una capa externa de comunicación, exposición o circulación disciplinada.
\end{definition}

\begin{observation}
La salida pública no equivale por sí sola a demostración final absoluta. Su función es fijar cuándo una afirmación ya puede abandonar el circuito interno de curación sin perder sobriedad, trazabilidad ni control de alcance.
\end{observation}

\section{Condición mínima de salida}
\begin{definition}
Un claim puede considerarse apto para salida pública cuando concurren simultáneamente:
\begin{enumerate}
  \item trazabilidad suficientemente limpia hacia su demo o evidencia de origen;
  \item dependencia canónica visible;
  \item alcance declarado de forma sobria;
  \item condiciones mínimas de validez y de fallo legibles;
  \item compatibilidad con auditoría mínima y ausencia de fragilidad crítica no declarada.
\end{enumerate}
\end{definition}

\begin{observation}
No se exige todavía clausura total del corpus. Se exige solo que la afirmación ya no dependa de opacidad estructural, exageración retórica o reserva implícita no expresada.
\end{observation}

\section{Salida pública y tabla madre}
\begin{definition}
La salida pública debe apoyarse en la tabla madre de demos y claims, de modo que toda afirmación proyectada conserve identificación, origen, estado editorial y relación visible con el resto del mapa metodológico.
\end{definition}

\begin{observation}
Esta cláusula evita que el claim público aparezca como enunciado aislado o desprendido de su genealogía interna.
\end{observation}

\section{Salida pública y sobriedad}
\begin{definition}
Toda salida pública debe distinguir expresamente entre:
\begin{enumerate}
  \item lo que la evidencia sostiene de manera efectiva;
  \item lo que la formulación interpreta de modo disciplinado;
  \item lo que permanece abierto como hipótesis, extensión o reserva.
\end{enumerate}
\end{definition}

\begin{observation}
La proyección externa no autoriza a borrar la diferencia entre soporte efectivo y proyección. Al contrario, exige volverla más visible.
\end{observation}

\section{Condición mínima de no salida}
\begin{definition}
Un claim no debe salir todavía al plano público cuando su formulación depende de reservas internas no explicitadas, su alcance sigue siendo inestable o su fuerza comunicable excede la fuerza metodológica realmente alcanzada.
\end{definition}

\begin{observation}
La no salida pública no degrada el valor del claim. Puede significar solo que todavía pertenece al circuito interno de curación, auditoría o reformulación.
\end{observation}

\section{Cláusula de alcance}
\begin{observation}
El criterio de salida pública de claims fija la frontera mínima entre curación interna y proyección externa disciplinada. No constituye todavía una política final de publicación integral ni una teoría exhaustiva de recepción pública del corpus.
\end{observation}

